\uuid{Zg0B}
\titre{Matrice d'une application linéaire}
\auteur{hanani ; mzari}
\organisation{exo7}
\licence{CC BY-NC-SA 4.0}
\datecreate{2020-01-20}
\niveau{L1}
\module{Algèbre}
\chapitre{Matrice}
\sousChapitre{Matrice et application linéaire}
\difficulte{1}
\theme{Matrice d'une application linéaire}
\exo7qcmid{lille-2.matrice-application.005}

\contenu{
\begin{question}
\qclassification{108.03}
\qtags{motcle=Matrice d'une application linéaire}
On considère $\Rr^3$ muni de sa base canonique et $f$ l'application linéaire définie par :
$$\begin{array}{rccc}f:&\Rr^3 &\to& \Rr^3\\
& (x,y,z)&\to &(x+y,x-z,y+z).  \end{array}$$
Quelles sont les assertions vraies ?
\begin{answers}
\good{La matrice de $f$ dans la base canonique est : $
\left(\begin{array}{rcc}
1&1&0\\
1&0&-1\\
0&1&1\\
\end{array}\right)$}.
\bad{La matrice de $f$ dans la base canonique est : $
\left(\begin{array}{rcc}
1&1&0\\
1&0&1\\
0&-1&1\\
\end{array}\right)$}.
\good{Le rang de $f$ est 2}.
\bad{Le rang de $f$ est 3}.
\end{answers}
\begin{explanations} Soit ${\cal {B}}=\{e_1,e_2,e_3\}$ la base canonique de $\Rr^3$. La matrice de $f$ dans la base  ${\cal {B}}$ est la matrice  dont la $j$ième colonne est constituée des coordonnées de  $f(e_j)$ dans la base ${\cal {B}}$. Cette matrice est : $
\left(\begin{array}{rcc}
1&1&0\\
1&0&-1\\
0&1&1\\ \end{array}\right)$.
\vskip0mm
On vérifie que $f(e_3)=f(e_2)-f(e_1)$ et que $ f(e_1)$ et $ f(e_2)$ ne sont pas colinéaires, donc  $ \{f(e_1), f(e_2)\}$
est une base de $\Im f$ et donc $\mbox{rg} (f)= \dim \Im f = 2$.
\end{explanations}
\end{question}
}
